% =====================================================================
%  Bölüm 7 — Doğal Sayılar ve Tümevarım
% =====================================================================
\chapter{Doğal Sayılar ve Tümevarım}\label{bol:tumevarim}

\begin{bolumhedefleri}
\begin{itemize}
  \item Peano aksiyomlarını ve onların \Lean{}'de neden \emph{teorem}
        olduğunu,
  \item Doğal sayıları kendi tümevarımsal tipimizle sıfırdan tanımlamayı,
  \item Toplama ve çarpmayı \emph{özyinelemeyle} tanımlamayı ve bu
        tanımların nasıl hesaplandığını,
  \item Tümevarım ilkesini, \leantac{induction} taktiğini ve tümevarım
        hipotezini kullanmayı,
  \item Toplamanın değişme ve birleşme özelliklerini ilk ilkelerden
        ispatlamayı,
  \item Gauss toplam formülünü, güçlü tümevarımı ve pul problemini,
  \item \leantac{simp}, \leantac{omega} ve \leantac{decide} otomasyonunu
        ve sınırlarını.
\end{itemize}
\end{bolumhedefleri}

Şimdiye kadar mantığın \emph{genel} yasalarıyla uğraştık: Her $P$, $Q$
ya da her yüklem için geçerli olan yasalar. Bu bölümde ilk kez belirli bir
matematiksel yapıyı, doğal sayıları ele alıyoruz. Doğal sayılar, mantıkla
matematiğin buluştuğu ilk yerdir. Burada öğreneceğimiz iki fikir,
\emph{özyinelemeli tanım} ve \emph{tümevarımla ispat}, kitabın geri
kalanının neredeyse her sayfasında karşımıza çıkacak.

% ---------------------------------------------------------------------
\section{Peano aksiyomları}\label{sec:peano}

Doğal sayılar nedir? Okulda $0, 1, 2, 3, \ldots$ diye sayarız; ama bu
``$\ldots$'' ne anlama gelir? On dokuzuncu yüzyılın sonunda Richard
Dedekind ve Giuseppe Peano bu soruya kesin bir cevap verdiler
\citep{dedekind1888, peano1889}. Doğal sayıların bütün özellikleri,
birkaç basit ilkeden türetilebilir (\cref{tab:peano}).

\begin{table}[htbp]
\centering
\renewcommand{\arraystretch}{1.5}
\rowcolors{2}{tablosatir}{white}
\begin{tabular}{>{\centering\arraybackslash}p{0.6cm}
                >{\raggedright\arraybackslash}p{6.2cm}
                >{\raggedright\arraybackslash}p{4.8cm}}
\rowcolor{tabloust}
\tablobaslik{} & \tablobaslik{Aksiyom} & \tablobaslik{Lean'de} \\
P1 & $0$ bir doğal sayıdır. & kurucu \lc{sifir} \\
P2 & Her $n$ doğal sayısının bir \emph{ardılı} $S(n)$ vardır. & kurucu \lc{ard} \\
P3 & $0$ hiçbir sayının ardılı değildir: $S(n) \neq 0$. & teorem \lc{sifir\_ard\_degil} \\
P4 & Ardıl fonksiyonu birebirdir: $S(n) = S(m) \Rightarrow n = m$. & teorem \lc{ard\_birebir} \\
P5 & \textbf{Tümevarım}: $0$'ı içeren ve ardıl altında kapalı olan her
     özellik, bütün doğal sayılarda geçerlidir. & teorem \lc{tumevarim} \\
\end{tabular}
\caption{Peano aksiyomları ve \Lean{}'deki karşılıkları. İlk ikisi
tanımın kendisidir; geri kalan üçü, tanımdan \emph{türetilen}
teoremlerdir.}
\label{tab:peano}
\end{table}

\begin{tarihce}[Tarihçe: Dedekind ve Peano]
Dedekind 1888'de yayımladığı \emph{Was sind und was sollen die Zahlen?}
(Sayılar nedir ve ne olmalıdır?) adlı kitabında doğal sayıları
kümeler kuramı içinde tanımladı ve özyinelemeli tanımların neden
anlamlı olduğunu ispatladı \citep{dedekind1888}. Peano bir yıl sonra
aynı fikirleri, kendi geliştirdiği sembolik dille ve aksiyomlar hâlinde
sundu \citep{peano1889}. Bugün kullandığımız $\in$ sembolü ve pek çok
mantık gösterimi Peano'dan gelir. Aksiyomlara Peano'nun adı verilir,
ama Peano kendisi Dedekind'in önceliğini kabul etmiştir.
\end{tarihce}

% ---------------------------------------------------------------------
\section{Doğal sayıları sıfırdan tanımlamak}\label{sec:dogal-tanim}

\Lean{}'in kendi doğal sayıları (\lc{Nat}) vardır; ama ne olduklarını
anlamanın en iyi yolu, onları kendimiz tanımlamaktır:

\leandosya[Kendi doğal sayılarımız]{Bolum07}{dogal_tanim}

\noindent \leankw{inductive} anahtar kelimesi yeni bir tip tanımlar ve
bu tipin elemanlarının \emph{nasıl kurulabileceğini} söyler. \lc{Dogal}
tipinin iki kurucusu vardır: \lc{sifir} ve \lc{ard}. Bu tipin her elemanı
bu iki kurucunun sonlu sayıda uygulanmasıyla elde edilir; başka hiçbir
eleman yoktur (\cref{fig:dogal-zincir}).

\begin{figure}[htbp]
\centering
\begin{tikzpicture}[
    sayi/.style={draw=anarenk, fill=tanimbg, thick, rounded corners=4pt,
                 font=\leanfont\scriptsize, minimum height=0.8cm, inner sep=4pt},
    ok/.style={-{Stealth[length=2.4mm]}, thick, leankw},
  ]
  \node[sayi] (s0) at (0,0) {sifir};
  \node[sayi] (s1) at (2.5,0) {ard sifir};
  \node[sayi] (s2) at (5.6,0) {ard (ard sifir)};
  \node[sayi] (s3) at (9.6,0) {ard (ard (ard sifir))};
  \node (s4) at (12.5,0) {$\cdots$};
  \foreach \a/\b in {s0/s1, s1/s2, s2/s3, s3/s4}
    \draw[ok] (\a) -- node[above, font=\leanfont\scriptsize\bfseries, text=leankw]{ard} (\b);
  \foreach \n/\v in {s0/0, s1/1, s2/2, s3/3}
    \node[below=4pt of \n, font=\baslikfont\small\bfseries, text=anarenk] {$\v$};
\end{tikzpicture}
\caption{\lc{Dogal} tipinin elemanları. Her sayı, \lc{sifir}'dan başlayıp
\lc{ard}'ı yeterince kez uygulayarak elde edilir.}
\label{fig:dogal-zincir}
\end{figure}

Tümevarımsal bir tanımın iki önemli özelliği vardır:
\begin{enumerate}[label=(\roman*)]
  \item \textbf{Kurucular birbirinden farklıdır}: \lc{sifir} ile
        \lc{ard n} asla eşit olamaz.
  \item \textbf{Kurucular birebirdir}: \lc{ard n = ard m} ise \lc{n = m}.
\end{enumerate}

\noindent Bu iki özellik \Lean{}'in tip kuramına yerleşiktir. Dolayısıyla
Peano'nun P3 ve P4 aksiyomları, bizim için \emph{ispatlanabilir
teoremlerdir}:

\leandosya[Peano aksiyomları teorem olarak]{Bolum07}{peano_teorem}

\begin{onemlifikir}[Önemli Fikir: Aksiyom değil, tanım]
Klasik matematikte Peano aksiyomları \emph{varsayılır}. Tip kuramında
ise tek yaptığımız \lc{Dogal} tipini tanımlamaktır; aksiyomların hepsi
bu tanımdan teorem olarak çıkar. Yeni bir varsayım eklemediğimiz için
bu yaklaşım, sistemin tutarlılığı açısından daha güvenlidir. Bu yöntemin
genel hâlini, yani her türlü veri yapısını \leankw{inductive} ile
tanımlamayı, \cref{bol:tumevarimsal-tipler}'da göreceğiz.
\end{onemlifikir}

% ---------------------------------------------------------------------
\section{Özyineleme: toplama ve çarpma}\label{sec:ozyineleme}

Doğal sayılar üzerinde fonksiyonları \emph{özyinelemeyle} tanımlarız:
Fonksiyonun \lc{sifir}'daki değerini ve \lc{ard n}'deki değerini
\lc{n}'deki değeri cinsinden veririz.

\leandosya[Toplama]{Bolum07}{topla_tanim}

\noindent Bu tanım, ikinci argüman üzerinde özyinelemelidir. ``$n + 0 = n$''
ve ``$n + (m+1) = (n+m)+1$'' kurallarını söyler. Bir toplamın nasıl
hesaplandığını adım adım izleyelim (\cref{fig:hesaplama-izi}).

\begin{figure}[htbp]
\centering
\begin{tikzpicture}[
    adim/.style={font=\leanfont\small, anchor=west},
    kural/.style={font=\baslikfont\scriptsize, text=leantac, anchor=west},
  ]
  \node[adim] at (0,0)    {topla iki (ard (ard sifir))};
  \node[adim] at (0,-0.75) {= ard (topla iki (ard sifir))};
  \node[adim] at (0,-1.5) {= ard (ard (topla iki sifir))};
  \node[adim] at (0,-2.25) {= ard (ard iki)};
  \node[adim] at (0,-3.0) {= ard (ard (ard (ard sifir)))};
  \node[kural] at (7.4,-0.75) {2. kural};
  \node[kural] at (7.4,-1.5)  {2. kural};
  \node[kural] at (7.4,-2.25) {1. kural};
  \node[kural] at (7.4,-3.0)  {\lc{iki}'nin tanımı};
  \node[font=\baslikfont\bfseries\small, text=anarenk, anchor=west] at (9.4,-1.5) {$2 + 2 = 4$};
  \draw[anarenk, thick, decorate, decoration={brace, amplitude=6pt}]
       (9.2,0.2) -- (9.2,-3.2);
\end{tikzpicture}
\caption{$2 + 2$'nin hesaplanması. Her adımda tanımdaki iki kuraldan biri
uygulanır; ikinci argüman her adımda küçülür ve hesaplama sona erer.}
\label{fig:hesaplama-izi}
\end{figure}

Çarpma da toplamaya dayanarak aynı biçimde tanımlanır:

\leandosya[Çarpma]{Bolum07}{carp_tanim}

\noindent Tanımlarımızı denemek için, kendi sayılarımızı \Lean{}'in
\lc{Nat}'ına çeviren bir fonksiyon yazıp \lc{\textcolor{leancmd}{\bfseries\#eval}}
ile hesaplatabiliriz:

\leandosya[Hesaplama]{Bolum07}{sayiya}

\begin{dikkat}[Dikkat: Özyineleme her zaman sona ermelidir]
\Lean{}, özyinelemeli her tanımın sona erdiğini denetler. Yukarıdaki
tanımlarda her özyinelemeli çağrı \emph{daha küçük} bir argümanla
(\lc{ard m} yerine \lc{m} ile) yapıldığı için bu denetim kendiliğinden
geçer. Buna \emph{yapısal özyineleme} denir. Sona ermeyen bir tanım
(örneğin \lc{f n := f n + 1}) kabul edilseydi, \lc{f 0 = f 0 + 1}
eşitliğinden \lc{False} türetilebilirdi. Sona erme denetimi, mantığın
tutarlılığının bir parçasıdır.
\end{dikkat}

% ---------------------------------------------------------------------
\section{Tümevarım ilkesi}\label{sec:tumevarim-ilkesi}

\begin{center}
\begin{prooftree}
  \hypo{P(0)}
  \hypo{[P(n)]^{1}}
  \ellipsis{}{P(n+1)}
  \infer2[Tümevarım$^{1,*}$]{\forall n,\; P(n)}
\end{prooftree}
\qquad
{\small \leantac{induction n with \ldots}}
\end{center}

\noindent {\small $^{*}$\,$n$, sağdaki türetimin diğer açık
varsayımlarında geçmemelidir.}

Tümevarım ilkesi, dominoların düşmesine benzetilir
(\cref{fig:domino}). Birinci domino düşüyorsa (\emph{taban adımı}) ve
her domino düştüğünde bir sonrakini de düşürüyorsa (\emph{tümevarım
adımı}), bütün dominolar düşer.

\begin{figure}[htbp]
\centering
\begin{tikzpicture}
  \draw[black!30, thick] (-0.9,0) -- (11.4,0);
  % düşmüş dominolar (sağ alt köşeleri etrafında döndürülmüş)
  \foreach \i in {0,1,2,3} {
    \draw[bittiframe, thick, fill=bittibg, rotate around={-58:({1.25*\i+0.2},0)}]
         ({1.25*\i-0.2},0) rectangle ({1.25*\i+0.2},1.5);
    \node[font=\scriptsize, text=bittiframe] at ({1.25*\i},-0.35) {$P(\i)$};
  }
  % ayakta olanlar
  \foreach \i in {4,5,6,7} {
    \draw[anarenk, thick, fill=tanimbg] ({1.25*\i-0.2},0) rectangle ({1.25*\i+0.2},1.5);
    \node[font=\scriptsize, text=anarenk] at ({1.25*\i},-0.35) {$P(\i)$};
  }
  \node at (10.6,0.7) {$\cdots$};
  \draw[-{Stealth[length=2.6mm]}, very thick, vurgu] (-0.8,2.0) -- (-0.15,1.2);
  \node[font=\baslikfont\scriptsize\bfseries, text=vurgu, anchor=south] at (-0.8,2.0) {taban: $P(0)$};
  \draw[-{Stealth[length=2.6mm]}, very thick, leantac] (4.0,2.3) to[bend left=20] (4.75,1.75);
  \node[font=\baslikfont\scriptsize\bfseries, text=leantac, anchor=south] at (4.0,2.3) {adım: $P(n) \to P(n+1)$};
\end{tikzpicture}
\caption{Tümevarım: İlk domino düşer ve her düşen domino bir sonrakini
düşürür; o hâlde hepsi düşer.}
\label{fig:domino}
\end{figure}

\Lean{}'de tümevarım ilkesi de bir \emph{teoremdir}. \leankw{inductive}
ile her tip tanımladığımızda, \Lean{} o tip için bir \emph{özyineleyici}
(İng.\ \emph{recursor}) üretir:

\leandosya[Özyineleyici]{Bolum07}{rec_check}

\noindent Bu tip, tam olarak tümevarım ilkesidir: \lc{motive} ispatlamak
istediğimiz özellik, ikinci argüman taban adımı, üçüncü argüman
tümevarım adımıdır. Aynı özyineleyici hem \emph{fonksiyon tanımlamak}
(\lc{motive} bir \lc{Type} ise) hem de \emph{ispat yapmak} (\lc{motive}
bir \lc{Prop} ise) için kullanılır. Curry--Howard'ın bir başka
yüzü: \textbf{tümevarımla ispat, özyinelemeli bir fonksiyondur}.

\leandosya[Tümevarım ilkesi]{Bolum07}{tumevarim}

\noindent \leantac{induction n with} satırı iki hedef açar. İkinci
hedefte bağlama \emph{tümevarım hipotezi} \lc{ih : P m} eklenir:

\begin{hedefgecisi}
\begin{hedefdurumu}[\texttt{sifir} durumu]
P : Dogal → Prop
h0 : P sifir
hs : ∀ (n : Dogal), P n → P n.ard
⊢ P sifir
\end{hedefdurumu}
\sonra
\begin{hedefdurumu}[\texttt{ard} durumu]
P : Dogal → Prop
h0 : P sifir
hs : ∀ (n : Dogal), P n → P n.ard
m : Dogal
ih : P m
⊢ P m.ard
\end{hedefdurumu}
\end{hedefgecisi}

% ---------------------------------------------------------------------
\section{İlk tümevarım ispatları}\label{sec:ilk-tumevarim}

\subsection{Bir asimetri}

Toplamanın tanımı ``$n + 0 = n$'' der. Bu eşitlik tanım gereği doğrudur;
\lc{rfl} yeter:

\leandosya[$n + 0 = n$]{Bolum07}{topla_sifir}

\noindent Ya ``$0 + n = n$''? Matematik açısından aynı derecede apaçık
görünür, ama \Lean{} için hiç de öyle değildir: Tanım ikinci argümana
bakar ve \lc{n} bilinmeyen bir sayı olduğu için hiçbir kural
uygulanamaz. Bu eşitlik ancak \emph{tümevarımla} ispatlanabilir:

\leandosya[$0 + n = n$]{Bolum07}{sifir_topla}

\noindent \leantac{show} adımı, hedefteki \lc{topla sifir (ard m)}
ifadesini tanımın ikinci kuralıyla bir adım açar. Ardından tümevarım
hipotezi \lc{ih : topla sifir m = m} ile yeniden yazmak yeterlidir.

\begin{onemlifikir}[Önemli Fikir: Tanımsal ve önermesel eşitlik]
\lc{topla n sifir = n} \emph{tanımsal} bir eşitliktir: \Lean{} iki tarafı
hesaplayarak aynı olduklarını görür ve \lc{rfl} kabul edilir.
\lc{topla sifir n = n} ise \emph{önermesel} bir eşitliktir: Doğrudur,
ama bir ispat gerektirir. Bu ayrım tip kuramının en önemli
kavramlarından biridir ve \cref{bol:bagimli}'te ayrıntılı olarak
incelenecek.
\end{onemlifikir}

Aynı ispatı terim modunda, \emph{özyinelemeli bir fonksiyon} olarak da
yazabiliriz. Kendini daha küçük bir argümanla çağıran bu ``ispat'',
tümevarımın kendisidir:

\leandosya[$0 + n = n$, özyinelemeli ispat]{Bolum07}{sifir_topla_terim}

\subsection{Değişme ve birleşme}

Toplamanın değişme özelliğini ($n + m = m + n$) ispatlamak için önce bir
yardımcı önsav gerekir:

\leandosya[Ardıl ile toplama]{Bolum07}{ard_topla}
\leandosya[Değişme özelliği]{Bolum07}{topla_degisme}
\leandosya[Birleşme özelliği]{Bolum07}{topla_birlesme}

\begin{figure}[htbp]
\centering
\begin{tikzpicture}[
    onsav/.style={draw=#1, fill=#1!8, thick, rounded corners=4pt,
                  font=\leanfont\scriptsize, align=center, inner sep=5pt,
                  minimum width=2.8cm},
    ok/.style={-{Stealth[length=2.4mm]}, thick, black!45},
  ]
  \node[onsav=anarenk] (t0) at (0,0)    {topla\_sifir\\{\rmfamily\small $n + 0 = n$}\\{\rmfamily\scriptsize tanımdan}};
  \node[onsav=ikincil] (s0) at (4.4,1.2)  {sifir\_topla\\{\rmfamily\small $0 + n = n$}};
  \node[onsav=ikincil] (at) at (4.4,-1.2) {ard\_topla\\{\rmfamily\small $(n{+}1) + m = (n{+}m){+}1$}};
  \node[onsav=teoremframe] (dg) at (9.0,0) {topla\_degisme\\{\rmfamily\small $n + m = m + n$}};
  \node[onsav=teoremframe] (bl) at (9.0,-2.5) {topla\_birlesme\\{\rmfamily\small $(a{+}b){+}c = a{+}(b{+}c)$}};
  \draw[ok] (t0) -- (s0);
  \draw[ok] (s0) -- (dg);
  \draw[ok] (at) -- (dg);
  \node[font=\baslikfont\scriptsize, text=black!60, align=center] at (0,-2.3) {her ok:\\``şu önsav kullanılır''};
\end{tikzpicture}
\caption{Toplamanın temel özellikleri arasındaki bağımlılıklar. Birleşme
özelliği hiçbir önsava ihtiyaç duymaz; değişme özelliği ise iki önsav
gerektirir.}
\label{fig:onsav-grafi}
\end{figure}

\Cref{fig:onsav-grafi}, ispatlar arasındaki bağımlılıkları gösteriyor.
Bu küçük grafik, büyük formalizasyon projelerinin nasıl planlandığının
bir minyatürüdür: Önce hangi önsavların gerektiği belirlenir, sonra
bunlar tek tek ispatlanır (\cref{bol:proje}).

\begin{dikkat}[Dikkat: Hangi değişken üzerinde tümevarım?]
Toplama ikinci argümanda özyinelemeli olduğu için, tümevarımı da
genellikle ikinci argüman üzerinde yapmak gerekir. \lc{topla\_degisme}
ispatında tümevarımı \lc{n} üzerinde yapmayı deneyin: Tanım
açılamayacağı için takılırsınız. Doğru değişkeni seçmek, tümevarım
ispatlarındaki en önemli karardır ve çoğunlukla \emph{tanımın} hangi
argümana baktığına göre verilir.
\end{dikkat}

% ---------------------------------------------------------------------
\section{Lean'in doğal sayıları}\label{sec:nat}

\Lean{}'in yerleşik \lc{Nat} tipi, bizim \lc{Dogal} tipimizle tam olarak
aynı biçimde tanımlanmıştır:

\leandosya[\texttt{Nat}'ın tanımı]{Bolum07}{nat_print}

\noindent Tek fark, \Lean{}'in \lc{Nat} için \lc{0}, \lc{1}, \lc{2}
gibi rakamları, \lc{n + 1} gösterimini ve büyük sayılarla verimli
hesaplamayı desteklemesidir. \leantac{induction} taktiği \lc{Nat} üzerinde
de aynı biçimde çalışır; durumlar \lc{zero} ve \lc{succ} adını taşır:

\leandosya[$0 + n = n$, \texttt{Nat} ile]{Bolum07}{nat_sifir_arti}

\noindent Bu ispatta \Lean{}'in kütüphanesindeki \lc{Nat.add\_assoc}
(birleşme) teoremini ters yönde kullandık. Bir sonraki örnekte ise
aritmetik işi \leantac{omega}'ya bırakıyoruz:

\leandosya[İki katı]{Bolum07}{iki_kat}

\subsection{Gauss'un toplamı}

Anlatılan hikâyeye göre küçük Carl Friedrich Gauss'a öğretmeni
1'den 100'e kadar olan sayıları topla demiş, Gauss da birkaç saniye
içinde 5050 cevabını vermiş. Gauss'un fikri, toplamı iki kez yazıp
terimleri eşlemekti (\cref{fig:gauss}).

\begin{figure}[htbp]
\centering
\begin{tikzpicture}[
    top/.style={circle, fill=#1, inner sep=0pt, minimum size=0.36cm},
  ]
  \foreach \satir in {1,...,5} {
    \foreach \sutun in {1,...,6} {
      \pgfmathtruncatemacro{\sinir}{\satir}
      \ifnum\sutun>\sinir
        \node[top=vurgu!55] at ({0.48*\sutun},{-0.48*\satir}) {};
      \else
        \node[top=anarenk] at ({0.48*\sutun},{-0.48*\satir}) {};
      \fi
    }
  }
  \draw[decorate, decoration={brace, amplitude=5pt}, thick] (0.25,-0.2) -- (3.12,-0.2)
       node[midway, above=6pt, font=\small] {$n + 1 = 6$};
  \draw[decorate, decoration={brace, amplitude=5pt, mirror}, thick] (0.18,-0.28) -- (0.18,-2.68)
       node[midway, left=6pt, font=\small] {$n = 5$};
  \node[font=\small, align=left, anchor=west] at (3.8,-1.0)
       {\textcolor{anarenk}{$\bullet$} $1 + 2 + \cdots + n$\\[2pt]
        \textcolor{vurgu!80}{$\bullet$} $n + \cdots + 2 + 1$};
  \node[font=\small, anchor=west] at (3.8,-2.2)
       {$2 \cdot (1 + \cdots + n) = n(n+1)$};
\end{tikzpicture}
\caption{Gauss'un fikri: İki üçgen birleşince $n \times (n+1)$'lik bir
dikdörtgen oluşur.}
\label{fig:gauss}
\end{figure}

\leandosya[Toplam fonksiyonu]{Bolum07}{gauss_tanim}

\noindent \lc{\textcolor{leancmd}{\bfseries\#eval}} ve \leantac{decide},
formülü belirli değerler için denetler; ama \emph{her} $n$ için
doğruluğunu göstermez. Bunun için tümevarım gerekir:

\begin{teorem}{Gauss}{gauss}
Her $n$ doğal sayısı için $2 \cdot (1 + 2 + \cdots + n) = n(n+1)$.
\end{teorem}

\leandosya[Gauss toplamı]{Bolum07}{gauss}

\noindent İspatın tümevarım adımı, kâğıt üzerindeki hesabın
satır satır karşılığıdır:
\begin{align*}
2 \cdot \bigl(T(m) + (m+1)\bigr)
  &= 2 \cdot T(m) + 2(m+1)   && \text{dağılma} \\
  &= m(m+1) + 2(m+1)          && \text{tümevarım hipotezi} \\
  &= (m+2)(m+1)               && \text{ters dağılma} \\
  &= (m+1)(m+2).              && \text{değişme}
\end{align*}

% ---------------------------------------------------------------------
\section{Güçlü tümevarım}\label{sec:guclu}

Bazı ispatlarda $P(n+1)$'i göstermek için yalnızca $P(n)$ değil, daha
önceki bütün değerler gerekir. Örneğin bir sayıyı asal çarpanlarına
ayırırken $n$'nin çarpanlarına bakarız ve bunlar $n-1$'den çok daha
küçük olabilir. \emph{Güçlü tümevarım} bu durumlar içindir
(\cref{fig:guclu}).

\begin{figure}[htbp]
\centering
\begin{tikzpicture}[
    kutu/.style={draw=#1, fill=#1!10, thick, rounded corners=2pt,
                 minimum width=0.8cm, minimum height=0.8cm, font=\scriptsize},
  ]
  \node[font=\baslikfont\small\bfseries, text=anarenk, anchor=west] at (-0.5,1.1) {Sıradan tümevarım};
  \foreach \i in {0,...,4} \node[kutu=gri] (a\i) at ({1.0*\i},0) {$\i$};
  \node[kutu=leantac] (a5) at (5,0) {$5$};
  \draw[-{Stealth[length=2.2mm]}, thick, anarenk] (a4) to[bend left=40] (a5);
  \node[font=\scriptsize, text=black!60, anchor=west] at (5.6,0) {$P(4) \Rightarrow P(5)$};

  \node[font=\baslikfont\small\bfseries, text=ikincil, anchor=west] at (-0.5,-1.1) {Güçlü tümevarım};
  \foreach \i in {0,...,4} \node[kutu=ikincil] (b\i) at ({1.0*\i},-2.2) {$\i$};
  \node[kutu=leantac] (b5) at (5,-2.2) {$5$};
  \foreach \i in {0,...,4} \draw[-{Stealth[length=2.2mm]}, thick, ikincil] (b\i) to[bend right=35] (b5);
  \node[font=\scriptsize, text=black!60, anchor=west] at (5.6,-2.2) {$P(0), \ldots, P(4) \Rightarrow P(5)$};
\end{tikzpicture}
\caption{Sıradan tümevarımda yalnızca bir önceki değer, güçlü
tümevarımda ise bütün önceki değerler kullanılabilir.}
\label{fig:guclu}
\end{figure}

\begin{teorem}{Güçlü tümevarım}{guclu}
$P$ doğal sayılar üzerinde bir yüklem olsun. Her $n$ için ``$m < n$ olan
her $m$ için $P(m)$'' varsayımından $P(n)$ çıkıyorsa, her $n$ için $P(n)$
doğrudur.
\end{teorem}

\noindent Güçlü tümevarımı ayrı bir aksiyom olarak eklememize gerek
yoktur; sıradan tümevarımdan türetilebilir. Hile, daha güçlü bir önermeyi
(``$n$'den küçük her $m$ için $P(m)$'') sıradan tümevarımla ispatlamaktır:

\leandosya[Güçlü tümevarım]{Bolum07}{guclu}

\noindent İspatın ortasındaki \leantac{have} adımına dikkat edin:
\emph{İspatlanacak önermeyi güçlendirmek}, tümevarım ispatlarındaki en
güçlü tekniklerden biridir. Doğrudan \lc{P n}'yi ispatlamaya çalışınca
tümevarım hipotezi yetersiz kalır; güçlendirilmiş önermenin tümevarım
hipotezi ise tam ihtiyacımız olanı verir.

\begin{onemlifikir}[Önemli Fikir: İyi sıralama]
Güçlü tümevarım, \emph{iyi sıralama ilkesine} denktir: Doğal sayıların
boş olmayan her alt kümesinin bir en küçük elemanı vardır. ``En küçük
karşı örnek'' yöntemiyle yapılan bütün ispatlar (``varsayalım ki $P$'yi
sağlamayan sayılar var; en küçüğü $n$ olsun\ldots'') aslında güçlü
tümevarımdır. Bu yöntemi \cref{bol:tumevarimsal-tipler}'da, sona erme
ispatlarında yeniden kullanacağız.
\end{onemlifikir}

% ---------------------------------------------------------------------
\section{Otomasyon: simp, omega, decide}\label{sec:otomasyon-nat}

Doğal sayılarla ilgili ispatların büyük kısmı rutin hesaplardan
oluşur. \Lean{}'in bu işler için güçlü araçları vardır:

\leandosya[Otomasyon örnekleri]{Bolum07}{otomasyon}

\begin{table}[htbp]
\centering
\renewcommand{\arraystretch}{1.5}
\rowcolors{2}{tablosatir}{white}
\begin{tabularx}{\linewidth}{>{\raggedright\arraybackslash}p{2.0cm}
                             >{\raggedright\arraybackslash}X
                             >{\raggedright\arraybackslash}p{4.2cm}}
\rowcolor{tabloust}
\tablobaslik{Taktik} & \tablobaslik{Ne yapar?} & \tablobaslik{Ne yapamaz?} \\
\leantac{decide}  & Değişken içermeyen, karar verilebilir önermeleri hesaplayarak çözer. & Değişken içeren önermeler; çok büyük hesaplar. \\
\leantac{simp}    & Bilinen eşitliklerle (\lc{n + 0 = n} gibi) hedefi sadeleştirir. & Her zaman kapatmaz; sonuç tahmin edilemeyebilir. \\
\leantac{omega}   & Tam sayı ve doğal sayılar üzerinde \emph{doğrusal} aritmetiği (toplama, sabitle çarpma, $<$, $\leq$, $=$) tam olarak karar verir. & Değişkenlerin çarpımı ($x \cdot y$) içeren \emph{doğrusal olmayan} adımlar. \\
\end{tabularx}
\caption{Doğal sayılar için otomasyon araçları.}
\label{tab:otomasyon-nat}
\end{table}

\begin{dikkat}[Dikkat: omega doğrusal olmayan adımları atlayamaz]
Gauss ispatının son adımı olan $m(m+1) + 2(m+1) = (m+1)(m+2)$ eşitliği
\emph{doğrusal değildir}: Değişkenler birbiriyle çarpılıyor. \leantac{omega}
bu tür terimleri yalnızca ``atom'' olarak, yani bilinmeyen sayılar gibi
görür; $m \cdot (m+1)$ ile $(m+1) \cdot (m+2)$ arasındaki ilişkiyi
göremez. Bu yüzden Gauss ispatında dağılma ve değişme adımlarını kendimiz
yazdık. \Mathlib{}'teki \leantac{ring} taktiği bu tür polinom
eşitliklerini kendiliğinden çözer; onu \cref{bol:kumeler}'de tanıyacağız.
\end{dikkat}

\begin{tarihce}[Tarihçe: Presburger aritmetiği]
1929'da Mojżesz Presburger, yalnızca toplamayı içeren doğal sayı
aritmetiğinin \emph{karar verilebilir} ve \emph{tam} olduğunu gösterdi
\citep{presburger1929}: Bu dildeki her önermenin doğru mu yanlış mı
olduğu bir algoritmayla belirlenebilir. Çarpma eklendiğinde ise durum
tamamen değişir: Gödel iki yıl sonra, toplama ve çarpmayı birlikte içeren
aritmetiğin eksik olduğunu gösterecekti (\ciltiki{4}). \leantac{omega}
taktiği, William Pugh'un Presburger aritmetiği için geliştirdiği
``omega testi'' algoritmasının bir çeşitlemesine dayanır \citep{pugh1991}.
\end{tarihce}

% ---------------------------------------------------------------------
\section{Curry--Howard: tümevarım = özyineleme}\label{sec:curry-howard-4}

\begin{center}
\renewcommand{\arraystretch}{1.5}
\rowcolors{2}{tablosatir}{white}
\begin{tabular}{>{\raggedright\arraybackslash}p{4.2cm}
                >{\raggedright\arraybackslash}p{4.6cm}
                >{\raggedright\arraybackslash}p{3.0cm}}
\rowcolor{tabloust}
\tablobaslik{Mantık} & \tablobaslik{Programlama} & \tablobaslik{Lean} \\
tümevarımsal tanım            & veri tipi tanımı                & \leankw{inductive} \\
tümevarımla ispat             & özyinelemeli fonksiyon          & \leantac{induction}, örüntü eşleme \\
taban adımı                   & taban durum                     & \lc{| sifir => \ldots} \\
tümevarım hipotezi            & özyinelemeli çağrının sonucu    & \lc{ih}, \lc{f m} \\
tümevarım ilkesi              & özyineleyici                    & \lc{Dogal.rec} \\
\end{tabular}
\end{center}

% ---------------------------------------------------------------------
\section{Taktik özeti}

\begin{center}
\renewcommand{\arraystretch}{1.45}
\rowcolors{2}{tablosatir}{white}
\begin{tabularx}{\linewidth}{>{\raggedright\arraybackslash}p{4.2cm} X}
\rowcolor{tabloust}
\tablobaslik{Taktik} & \tablobaslik{Ne yapar?} \\
\leantac{induction n with | ... }  & \lc{n} üzerinde tümevarım; her kurucu için bir durum açar. \\
\leantac{injection h with h'}      & \lc{h : c a = c b}'den \lc{h' : a = b} çıkarır. \\
\leantac{cases h} (\lc{h : c₁ = c₂}) & Farklı kurucuların eşitliğinden çelişki çıkarır. \\
\leantac{show t}                   & Hedefi tanımsal olarak eşit bir biçime çevirir (tanımı açar). \\
\leantac{change t at h}            & Hipotezi tanımsal olarak eşit bir biçime çevirir. \\
\leantac{omega}                    & Doğrusal tam sayı aritmetiğini çözer. \\
\leantac{simp}                     & Hedefi bilinen eşitliklerle sadeleştirir. \\
\leantac{decide}                   & Değişken içermeyen önermeleri hesaplar. \\
\end{tabularx}
\end{center}

% ---------------------------------------------------------------------
\begin{arastirmanotu}
\textbf{Natural Number Game.} Kevin Buzzard ve arkadaşlarının hazırladığı
çevrimiçi \emph{Natural Number Game}, bu bölümdeki ispatları (ve çok
daha fazlasını) adım adım bir oyun olarak sunar \citep{nng}. Peano
aksiyomlarından başlayıp $\leq$ ilişkisine ve asal sayılara uzanır.
Bu bölümü bitiren okura özellikle öneririz.

\textbf{Tümevarımın gücü.} Peano aritmetiğinin (PA) gücü, tümevarım
aksiyomunun hangi önermeler için kabul edildiğine bağlıdır. Tümevarımı
yalnızca basit (sınırlı niceleyicili) önermelerle kısıtlayan zayıf
aritmetik sistemleri, karmaşıklık kuramıyla ilişkileri nedeniyle hâlâ
etkin bir araştırma alanıdır. Gentzen'in 1936'daki tutarlılık ispatı,
PA'nın tutarlılığının $\varepsilon_0$ ordinaline kadar
\emph{sonsuz ötesi tümevarımla} gösterilebileceğini ortaya koydu
\citep{gentzen1936}; bu konuya \ciltiki{4}'te döneceğiz.

\textbf{İleri okuma.} Tümevarım ve özyineleme için \citet{avigad2022};
\Lean{}'de özyinelemeli tanımlar ve sona erme için \citet{tpil4} ve
\citet{fpil}.
\end{arastirmanotu}

\begin{bolumozeti}
\begin{itemize}
  \item Doğal sayılar, \lc{sifir} ve \lc{ard} kurucularıyla tümevarımsal
        bir tip olarak tanımlanır. Peano aksiyomları bu tanımdan teorem
        olarak çıkar.
  \item Fonksiyonlar özyinelemeyle tanımlanır; \Lean{} her özyinelemenin
        sona erdiğini denetler.
  \item Tümevarım ilkesi, tipin özyineleyicisidir. \leantac{induction}
        taktiği taban ve adım durumlarını açar; adım durumunda
        tümevarım hipotezi \lc{ih} kullanılır.
  \item \lc{n + 0 = n} tanımsal, \lc{0 + n = n} önermeseldir ve tümevarım
        ister. Tümevarım genellikle tanımın özyinelemeli olduğu argüman
        üzerinde yapılır.
  \item Güçlü tümevarım sıradan tümevarımdan türetilir; anahtar teknik
        önermeyi güçlendirmektir.
  \item \leantac{omega} doğrusal aritmetiği, \leantac{decide} somut
        hesapları çözer. Doğrusal olmayan adımlar elle ya da \Mathlib{}'in
        \leantac{ring} taktiğiyle yapılır.
\end{itemize}
\end{bolumozeti}

\begin{alistirmalar}
\alistirma{1} \lc{topla n (ard sifir) = ard n}. Neden tümevarım
gerekmiyor? (\lc{a7\_1})

\alistirma{2} \lc{carp sifir n = sifir}. Hangi değişken üzerinde
tümevarım yapmalısınız? (\lc{a7\_2})

\alistirma{2} Toplamada sadeleştirme: \lc{topla a c = topla b c → a = b}.
İpucu: Tümevarım adımında \leantac{injection} kullanın. (\lc{a7\_3})

\alistirma{1} \lc{Nat} üzerinde, kütüphane teoremi kullanmadan:
\lc{0 * n = 0}. (\lc{a7\_4})

\alistirma{2} \lc{ucKat n = 3 * n}. (\lc{a7\_5})

\alistirma{2} Tek sayıların toplamı bir karedir:
$1 + 3 + \cdots + (2n-1) = n^2$. Önce \cref{fig:gauss}'e benzer bir
resimle neden doğru olduğunu gösterin, sonra \Lean{}'de ispatlayın.
(\lc{a7\_6})

\alistirma{2} \lc{n < 2 \textasciicircum{} n}. İpucu: \lc{Nat.pow\_succ}. (\lc{a7\_7})

\alistirma{3} \textbf{Pul problemi}: 8 kuruş ve üstündeki her tutar
yalnızca 3 ve 5 kuruşluk pullarla ödenebilir. Güçlü tümevarımla
ispatlayın. Hangi üç durum taban adımı olarak gerekli? (\lc{a7\_8})

\alistirma{2} (Kâğıt-kalem) \lc{carp iki uc} ifadesinin hesaplanmasını
\cref{fig:hesaplama-izi}'deki gibi adım adım yazın.

\alistirma{3} (Kâğıt-kalem) \lc{topla\_degisme}'yi \lc{n} üzerinde
tümevarımla ispatlamaya çalışın. Hangi adımda takılıyorsunuz?
Takılmamak için hangi önsavlar gerekir?

\alistirma{3} (Tartışma) Aşağıdaki ``ispatın'' hatası nedir?
``Bütün atlar aynı renktir. Tümevarımla: Bir attan oluşan her grupta
atlar aynı renktir. $n$ atlık gruplar için doğru olsun; $n+1$ atlık bir
grupta ilk $n$ at aynı renktir, son $n$ at da aynı renktir; iki grup
ortak atlar içerdiği için hepsi aynı renktir.''
\end{alistirmalar}
